% GENERATED FILE. Edit canonical lessons, then run npm run generate:books.
% canonical-source-hash: fnv1a64:503e0d52446ab136
% canonical-lessons: KA-C10-vaara, KA-S119-vowel-sign-oo, KA-S141-digit-two, KA-C10-checkpoint

\chapter{Days of the Week}
\label{ch:ka-days-of-the-week}
\hlchaptermodality{10}

\begin{chapteropening}
\textbf{By the end of this chapter:} \emph{I can name any day of the week in Kannada, including the Sunday that uses a different sun-name from Hindi's.}

Complete KA-C10-checkpoint, the chapter's closing checkpoint: say the seven days from sōmavāra to bhānuvāra and name the sun-name that differs from Hindi's, find \kn{◌ೋ} in \kn{ಸೋಮವಾರ} and read \kn{೨} as eraḍu, then write both from sound.
\end{chapteropening}

\section[\texorpdfstring{\kn{ಸೋಮವಾರ} \kn{ಮಂಗಳವಾರ} … (sōmavāra …)}{ಸೋಮವಾರ ಮಂಗಳವಾರ … (sōmavāra …)}]{\kn{ಸೋಮವಾರ} to \kn{ಭಾನುವಾರ} — the Sanskrit week}

\label{lesson:KA-C10-vaara}



\begin{quote}
Unlike Tamil and Malayalam, Kannada's week is built entirely from Sanskrit — the same recipe as Hindi's. But even here, one day picks a \textbf{different} Sanskrit word than Hindi does.
\end{quote}

\subsection*{You'll want to know: The pattern: {[}deity{]} + \kn{ವಾರ}, \textquotedblleft{}day\textquotedblright{}}
\noindent
\begin{tabularx}{\linewidth}{@{}>{\raggedright\arraybackslash}X>{\raggedright\arraybackslash}X>{\raggedright\arraybackslash}X@{}}
\toprule
\textbf{Kannada} & \textbf{deity/planet} & \textbf{day} \\
\midrule
\textbf{\kn{ಸೋಮವಾರ}} (\emph{Somavāra}) & \textbf{Moon} (Soma) & Monday \\
\textbf{\kn{ಮಂಗಳವಾರ}} (\emph{Maṅgaḷavāra}) & \textbf{Mars} (Maṅgala) & Tuesday \\
\textbf{\kn{ಬುಧವಾರ}} (\emph{Budhavāra}) & \textbf{Mercury} (Budha) & Wednesday \\
\textbf{\kn{ಗುರುವಾರ}} (\emph{Guruvāra}) & \textbf{Jupiter} (Guru/Bṛhaspati) & Thursday \\
\textbf{\kn{ಶುಕ್ರವಾರ}} (\emph{Śukravāra}) & \textbf{Venus} (Śukra) & Friday \\
\textbf{\kn{ಶನಿವಾರ}} (\emph{Śanivāra}) & \textbf{Saturn} (Śani) & Saturday \\
\textbf{\kn{ಭಾನುವಾರ}} (\emph{Bhānuvāra}) & the \textbf{Sun} (Bhānu, \textquotedblleft{}light, splendor\textquotedblright{}) & Sunday \\
\bottomrule
\end{tabularx}

Five of these are identical in shape to Hindi's — same deity, same \textbf{‑vāra/ ‑vār} ending. \textbf{Sunday} is the interesting exception: where Hindi says \textbf{\dv{रविवार}} (\emph{Ravivār}, from \emph{Ravi}, \textquotedblleft{}sun\textquotedblright{}), Kannada says \textbf{\kn{ಭಾನುವಾರ}} (\emph{Bhānuvāra}), built on a \textbf{different} Sanskrit word for the sun — \textbf{Bhānu}, \textquotedblleft{}light, splendor,\textquotedblright{} rather than \emph{Ravi}. Sanskrit has \textbf{many} names for the sun, and different languages downstream picked different ones for their week.

\subsection*{Your turn}
\begin{itemize}
\raggedright
  \item \emph{Say it:} \textquotedblleft{}Somavāra, Maṅgaḷavāra, Budhavāra, Guruvāra, Śukravāra, Śanivāra\textquotedblright{} — Monday through Saturday
  \item \emph{Say it:} the different one — \textquotedblleft{}Bhānuvāra,\textquotedblright{} Sunday
  \item \emph{Say it:} the contrast — Hindi says Ravivāra (Ravi), Kannada says Bhānuvāra (Bhānu) — two different Sanskrit sun-names
  \item \emph{Recall:} say \emph{ēḷu}
  \item \emph{Read it:} \textbf{\kn{ಎಂಟು}}
\end{itemize}

\begin{tcolorbox}[breakable,colback=teal!4,colframe=teal!35!black,title={Before you move on}]
What ending does every Kannada weekday share, and where's it from? (\textbf{‑\kn{ವಾರ}} \emph{-vāra}, Sanskrit \textquotedblleft{}day.\textquotedblright{}) Which day differs from Hindi's version, and how? (\textbf{Sunday} — Kannada's \emph{Bhānuvāra} (\textquotedblleft{}light\textquotedblright{}) uses a different Sanskrit sun-name than Hindi's \emph{Ravivār} (\textquotedblleft{}Ravi\textquotedblright{}).)
\end{tcolorbox}
\section[\texorpdfstring{\kn{ೋ} (ō)}{ೋ (ō)}]{\kn{◌ೋ} — one character, met inside words you already say}

\label{lesson:KA-S119-vowel-sign-oo}



\begin{quote}
Before the new one: \kn{ಡ} — what does it do?

One character this time. Just one — and you have been saying it for pages without knowing which mark on the page it was.
\end{quote}

\subsection*{Script you'll notice: \kn{◌ೋ}}
\textbf{\kn{◌ೋ}} — \emph{ō}.

It is a \textbf{vowel sign}. It is not a letter and never stands alone: it attaches to a consonant and \textbf{replaces} the \emph{a} built into it with \emph{ō}. Replaces, not adds — the \emph{a} is gone.

You already say these, and every one of them has \kn{◌ೋ} somewhere inside it:

\begin{itemize}
\raggedright
  \item \textbf{\kn{ಸಂತೋಷ}} \emph{santōṣa} — joy / pleased to meet you
  \item \textbf{\kn{ಹೋಗು}} \emph{hōgu} — to go (and \kn{ಬಾ}, come)
  \item \textbf{\kn{ಹೋಗಿ} \kn{ಬರುತ್ತೇನೆ}} \emph{hōgi baruttēne} — goodbye (lit. \textquotedblleft{}I'll go and come back\textquotedblright{})
  \item \textbf{\kn{ನಾಳೆ} \kn{ಸಿಗೋಣ}} \emph{nāḷe sigōṇa} — see you tomorrow
\end{itemize}

\subsection*{Writing: \kn{◌ೋ} — copy what you see}
Put your pen on \kn{◌ೋ} and follow its line. Copy the shape you can see — slowly, and larger than it is printed.

\begin{quote}
This book does not yet tell you \textbf{where to start the character or which way to travel}. That is a real thing, taught with real variation from school to school, and it is not written down here until it can be written down with a source. Copying what is in front of you needs no such source, and it is how the shape gets into your hand in the meantime.
\end{quote}

\subsection*{Your turn}
\begin{itemize}
\raggedright
  \item \emph{Look:} at these words, and find \kn{◌ೋ} in the ones that have it
\end{itemize}

\begin{quote}
\kn{ಸಂತೋಷ}  ·  \kn{ಹೋಗು}  ·  \kn{ನಮಸ್ಕಾರ}
\end{quote}

\begin{itemize}
\raggedright
  \item \emph{Trace it:} \kn{◌ೋ} three times, saying \emph{ō} as you finish each one
  \item \emph{Look:} back at any page of this chapter and find \kn{◌ೋ} once more
\end{itemize}

\begin{tcolorbox}[breakable,colback=teal!4,colframe=teal!35!black,title={Before you move on}]
Which character is this — \kn{◌ೋ}? What vowel does it put on the consonant it attaches to — and what does it take off? (\textbf{Puts \emph{ō} on; takes the built-in \emph{a} off.}) Name one word you already say that contains it.
\end{tcolorbox}
\section[\texorpdfstring{\kn{೨} (2)}{೨ (2)}]{\kn{೨} — the digit for eraḍu}

\label{lesson:KA-S141-digit-two}



\begin{quote}
Before the new shape: \textbf{\kn{೧}} — which number is it?

One digit this time, and you already know its name.
\end{quote}

\subsection*{Script you'll notice: \kn{೨}}
\textbf{\kn{೨}} — \textbf{2}, the digit for \textbf{\kn{ಎರಡು}} (\emph{eraḍu}), \textquotedblleft{}two\textquotedblright{}.

Each Dravidian script carries its own set of ten digits. Telugu, Malayalam and Tamil all have theirs; none of the four sets is borrowed from another.

\subsection*{Writing: \kn{೨} — follow the numbered strip}
\hlblockfigure{\detokenize{figures/KA-S141-digit-two-filmstrip.pdf}}{How \kn{೨} is written, stroke by stroke}

Follow the numbered strip: start where movement 1 begins, and let each panel show you the next movement before your pen makes it. Then write \kn{೨} yourself — slowly, and larger than it is printed.

\begin{quote}
The strip shows \textbf{where to start the character and which way to travel}, in one attested order. That is taught with real variation from school to school, so the strip names its source beneath it: treat it as a sound way in, not the only one.
\end{quote}

\subsection*{Your turn}
\begin{itemize}
\raggedright
  \item \emph{Look:} at the digits you have met, and name each one aloud
\end{itemize}

\begin{quote}
\kn{೧}  ·  \kn{೨}
\end{quote}

\begin{itemize}
\raggedright
  \item \emph{Say it:} \textquotedblleft{}eraḍu\textquotedblright{} as you point at \kn{೨}
\end{itemize}

\begin{tcolorbox}[breakable,colback=teal!4,colframe=teal!35!black,title={Before you move on}]
Which number is \textbf{\kn{೨}}? (\textbf{2} — \emph{eraḍu}, \kn{ಎರಡು}.) Name every digit you have met so far, in order.
\end{tcolorbox}
\section[Practice]{Chapter 10 checkpoint — the week, \kn{◌ೋ} and \kn{೨}}

\label{lesson:KA-C10-checkpoint}



\begin{quote}
Nothing new here. The week out loud first, then the two shapes this chapter put on the page.
\end{quote}

\subsection*{Your turn — the week}
\begin{itemize}
\raggedright
  \item \emph{Say it:} the seven days, Monday to Sunday
\end{itemize}

(\emph{Sōmavāra, maṅgaḷavāra, budhavāra, guruvāra, śukravāra, śanivāra, bhānuvāra}.)

\begin{itemize}
\raggedright
  \item What ending does every day share, and what does it mean? (\emph{-vāra}, \textquotedblleft{}day.\textquotedblright{})
  \item Which day takes a different Sanskrit sun-name from Hindi's, and what are the two names? (Sunday: Kannada \emph{bhānuvāra}, \textquotedblleft{}light\textquotedblright{}; Hindi \emph{ravivār}.)
\end{itemize}

\subsection*{Script — \kn{◌ೋ} and \kn{೨}}
\begin{itemize}
\raggedright
  \item \textbf{\kn{ಸೋಮವಾರ}} is Monday. Which mark in it gives its first consonant the \emph{ō}, and what does that mark take away? (\textbf{\kn{◌ೋ}}: it puts \emph{ō} on and takes the built-in \emph{a} off.)
  \item Which number is \textbf{\kn{೨}}, and which word does it stand for? (\textbf{2}, \emph{eraḍu}.)
\end{itemize}

\emph{Read it:} \textbf{\kn{ಸೋಮವಾರ}} aloud, and point to the \kn{◌ೋ} in it

\emph{Read it:} \textbf{\kn{೨}} aloud as \emph{eraḍu}

\subsection*{Writing — from sound}
\emph{Cover:} the Script block above, so no model stays in view

\emph{Write it:} the digit for \emph{eraḍu}, then the sign that turns a consonant's \emph{a} into \emph{ō}

\emph{Check:} against the key — \textbf{\kn{೨}}, \textbf{\kn{◌ೋ}}

\emph{Write it:} one repaired shape, then stop

\begin{tcolorbox}[breakable,colback=teal!4,colframe=teal!35!black,title={Before you move on}]
Say today's day in Kannada, then Sunday's. That is the chapter: the week by name, and two shapes that sit inside it on the page.

Next chapter: the colours.
\end{tcolorbox}
